
\paragraph{Ánh sáng trắng} Ánh sáng trắng là ánh sáng tập hợp của vô số ánh sáng (bức xạ) khác nhau trong đó có vùng nhìn thấy có màu biến thiên từ đỏ đến tím.
\paragraph{Đặc điểm về giao thoa ánh sáng trắng}
\begin{itemize}
\item Tại O là vân sáng trắng, hai bên là dải màu cầu vồng (tím trong, đỏ ngoài).
\item Bề rộng quang phổ là khoảng cách từ vân sáng đỏ đến vân sáng tím cùng bậc và cùng phía với vân trung tâm.
\begin{align*}
\boder{\Delta x_k=k(i_{\ct{đỏ}}-i_{\ct{tím}})=k\dfrac{\pbrace{\lambda_{\ct{đỏ}}-\lambda_{\ct{tím}}}D}{a}}
\end{align*}
\item Ánh sáng trắng gồm các bức xạ có bức xạ \bth{0,38\microm\leq \lambda\leq 0,76\microm}.
\end{itemize}
